\chapter{Критерий повторного открытия программы физического происхождения}
\label{ch:v9-physical-origin-reopening-criterion}

\section{Три открытых критерия}

После отделения conditional и physical ledger остаётся deficit
\begin{equation}
 \mathbf d_{\rm phys}=(1,1,1)^T,
 \label{eq:v9-reopening-deficit}
\end{equation}
соответствующий physical selection масштаба и связи, физический гессиан и
unconditional слепое безразмерное следствие.

Эти три program slots не являются тремя независимыми происхождениями.
Минимальная карта двух provenance-пакетов равна
\begin{equation}
 R_{\rm reopen}=
 \begin{pmatrix}
  1&0\\
  0&1\\
  0&1
 \end{pmatrix}.
 \label{eq:v9-reopening-map}
\end{equation}
Первый столбец означает physical origin общего $E_*$ и coupling $\chi$.
Второй означает physical origin логарифм определителя parent-term вместе с его
коэффициентом. После такого происхождения тот же term даёт и положительный
shape-Hessian, и blind normalized consequence.

\section{Минимальность reopening dossier}

Для совместного dossier
\begin{equation}
 \mathbf a_{\rm joint}=(1,1)^T,
 \qquad
 R_{\rm reopen}\mathbf a_{\rm joint}=\mathbf d_{\rm phys},
 \qquad
 \rank R_{\rm reopen}=2.
 \label{eq:v9-reopening-joint-closure}
\end{equation}
Ни один пакет отдельно не покрывает deficit:
\begin{equation}
 \rank\bigl(R_{\rm reopen}^{(1)}\mid\mathbf d_{\rm phys}\bigr)=2,
 \qquad
 \rank\bigl(R_{\rm reopen}^{(2)}\mid\mathbf d_{\rm phys}\bigr)=2.
 \label{eq:v9-reopening-single-package-no-go}
\end{equation}
Следовательно, необходимы два независимых physical-origin certificate,
хотя один из них закрывает сразу два program criteria.

Допустимый certificate должен быть source-free относительно искомого
числа: его microscopic carrier, measure/action и coefficient должны быть
заданы до чтения target ratios. Кроме того, оба пакета должны входить в
один ограниченный снизу common parent; две несвязанные подгонки не переоткрывают
программу Тома IX.

\section{Текущая граница}

Для augmented model условная availability равна
\begin{equation}
 \mathbf a_{\rm cond}=(1,1)^T,
 \qquad R_{\rm reopen}\mathbf a_{\rm cond}=(1,1,1)^T.
 \label{eq:v9-reopening-conditional-availability}
\end{equation}
Но admitted scale wells и логарифм определителя axiom не имеют physical-origin
certificate. Поэтому
\begin{equation}
 \mathbf a_{\rm phys}=(0,0)^T,
 \qquad R_{\rm reopen}\mathbf a_{\rm phys}=(0,0,0)^T,
 \label{eq:v9-reopening-physical-availability}
\end{equation}
и строгий residual остаётся
\begin{equation}
 \mathbf d_{\rm phys}-R_{\rm reopen}\mathbf a_{\rm phys}
 =(1,1,1)^T,
 \qquad \sum_i d_i^{\rm residual}=3>0.
 \label{eq:v9-reopening-residual}
\end{equation}

\begin{theorem}[Минимальный reopening criterion]
Физическая ветвь Тома IX может быть переоткрыта только совместным dossier
из двух independent provenance-пакетов: physical scale/coupling origin и
physical логарифм определителя-parent origin. Их incidence map имеет rank $2$ и покрывает
все три открытых program criteria.
\end{theorem}

\begin{nogocriterion}[Запрет фиктивного переоткрытия]
Один physical scale anchor без логарифм определителя origin либо один логарифм определителя mechanism
без physical $E_*,\chi$ не переоткрывает программу. Не допускаются также
содержащий целевую подстановку transformations, fitted coefficients и два несвязанных
parent-term, не происходящих из одного microscopic carrier.
\end{nogocriterion}

\begin{protocol}\begin{enumerate}
\item open physical criteria: \textbf{$3$};
\item minimal independent origin packages: \textbf{$2$};
\item conditional deficit coverage: \textbf{$3/3$};
\item physical deficit coverage: \textbf{$0/3$};
\item physical reopening packages: \textbf{$0/2$};
\item status freeze: \textbf{сохраняется};
\item следующий гейт: \textbf{common-origin carrier admission}.
\end{enumerate}\end{protocol}

Машинный сертификат сохранён в
\path{s2t_v9_axiom_augmented_physical_origin_reopening_criterion_gate_results.json}.